\chapter{Название}

\lesson{1}{от 25 фев 2025 12:45}{Начало}


\section{Введение}

\subsection{Содержание}

Часть 1
\begin{enumerate}
	\item Что такое ИИ
	\item Основные типы
	\item Компоненты ИИ
	\item Типы обучения
	\item  Примеры применения
\end{enumerate}

Часть 2
\begin{enumerate}
	\item Постановка задачи обучения
	\item Основные стандартные типы задач
\end{enumerate}

\begin{definition}[Искусственный интелект]
	\emph{Искусственный интелект} -- это имитация процессов человеческого интеллекта машинами, особенно компьютерными системами.
\end{definition}

\subsection{Какие же вообще сущестувуют задачи машинного обучения?}

Обычно выделяют следующие категории:

\begin{itemize}
	\item задачи классификации;
	\item задачи регрессии;
	\item задачи ранжирования.
\end{itemize}

\begin{definition}[Задача классификации]
	\emph{Задача классификации} -- это когда требуется входные данные отнести к тому или иному классу (например, изображения разделить на кошек и собак, или распознать произнесенное слово, или по медицинским данным выдать диагноз и т.п.)
\end{definition}

\begin{definition}[Задача регрессии]
	В \emph{задачах регрессии} обычно делают прогнозы в виде вещественных чисел на основе выходных данных.
\end{definition}

\begin{definition}[Задача ранжирования]
	\emph{Задачи ранжирования} подразумевают упорядочивание (по некоторому критерию) входного набора данных. Классический пример -- это поисковые системы, которые ранжируют поисковую выборку по релевантности для каждого конкретного пользователя.
\end{definition}

\begin{note}[Общая постановка задачи обучения по прецедентам]
	Начало: \\
	$X$ -- множество объектов; \\
	$Y$ -- множество ответов; \\
	$y: X \rightarrow Y$ -- неизвестная зависимость (target function). \\

	Дано: \\
	$\{x_1,\ldots,x_l\} \subset X$ -- обучающая выборка (training sample); \\
	$y_i = y(x_i), \ i = 1,\ldots ,l$ -- известные ответы. \\

	Найти: \\
	$a: X \rightarrow Y$ -- алгоритм, решающую функцию (decision function), приближающую $y$ на всем множестве $X$.
\end{note}

\begin{itemize}
	\item Формирование данных для обучения, зачастую, самый сложный и длительный процесс, требующий участие одного или группы людей.
	\item Исходные результаты измерений $\{x_i\}$ могут иметь сложную структуру (например, при кодировании текстов, звуков или изображений) и трудно интерпретируемыми самими алгоритмами.
\end{itemize}

\begin{note}[Признаковое описание объектов]
	$f_j: X \rightarrow D_j, \ j = \overline{1,n} $ -- признаки объектов (features).

	Типы признаков:
	\begin{itemize}
		\item $D_j = \{0,1\}$ -- бинарный признак $f_j$;
		\item $\abs{D_j} < \infty $ -- номинальный признак $f_j$;
		\item $\abs{D_j} < \infty , \ D_j$ упорядочено -- порядковый признак $f_j$;
		\item $D_j = \R$ -- количественный признак $f_j$.
	\end{itemize}

	Вектор $\big(f_1(x), \ldots , f_n(x)\big)$ -- признаковое описание объекта $x$.
\end{note}
